\documentclass[../linear-algebra.tex]{subfiles}
\begin{document}
\section{Eigenvectors/values}
\subsection{Definitions}
Eigenvalues and eigenvector apply only to square matrices.
\paragraph{Eigenvector} An eigenvector of an $n \times n$ matrix $A$ is a nonzero vector $\mathbf{x}$ such that the equation
$$A\mathbf{x}=\lambda \mathbf{x}$$
holds true for some scalar $\lambda$, called the \textbf{eigenvalue}.\\
By definition, an eigenvector must be nonzero to avoid the trivial and unimportant case where $A\mathbf{0}=\lambda\mathbf{0}$, which is always true for any matrix and scalar.\\
The valid vector $\mathbf{x}$ is referred to as the eigenvector corresponding to that specific $\lambda$.\\
Geoemetrically, when an eigenvector is multiplied by matrix $A$, its direction remains completely unchanged. It can be stretched or compressed.

\subsection{Eigenvalue can be 0}
If $\lambda = 0$, the core equation becomes $A\mathbf{x} = 0\mathbf{x}$, which simplifies to $A\mathbf{x} = \mathbf{0}$.\\
Because eigenvectors must be nonzero, this means $A\mathbf{x} = \mathbf{0}$ has a nontrivial solution.\\
Therefore, an $n \times n$ matrix $A$ having an eigenvalue of 0 means $A$ is not invertible.\\
\\
Geometrically, there is at least one non-zero direction in space (the eigenvector) that gets completely crushed down to the origin during the transformation.
So the matrix $A$ represents transformation that crushes space and loses a dimension.\\
An invertible matrix represents a transformation that can be perfectly reversed or undone. For matrix with eigenvalue of 0, we cannot reverse the operation of crushing the eigenvectors down to $\mathbf{0}$, because the unique starting points are lost.
Hence it ties in with matrix being singular.\\
\\
The eigenspace for an eigenvalue of $0$ is the same thing as the null space of the matrix.

\subsection{Finding eigenvectors from known eigenvalue}
Solving for nontrivial $\mathbf{x}$ to $A\mathbf{x} = \lambda\mathbf{x}$, rewrite it as
$$(A-\lambda I)\mathbf{x}=\mathbf{0}$$
The set of all possible solutions to $(A-\lambda I)\mathbf{x}=\mathbf{0}$ forms the null space of the matrix $(A-\lambda I)$. This set is a subspace of $\mathbb{R}^n$ and is formally called the eigenspace of $A$ corresponding to $\lambda$.
The eigenspace contains the zero vector and all the eigenvectors tied to that specific eigenvalue.
\paragraph{Example} Given a $3 \times 3$ matrix $A = \begin{bmatrix}4 & -1 & 6 \\ 2 & 1 & 6 \\ 2 & -1 & 8\end{bmatrix}$ with an eigenvalue of $\lambda=2$,
$$A - 2I = \begin{bmatrix} 4-2 & -1 & 6 \\ 2 & 1-2 & 6 \\ 2 & -1 & 8-2 \end{bmatrix} = \begin{bmatrix} 2 & -1 & 6 \\ 2 & -1 & 6 \\ 2 & -1 & 6 \end{bmatrix}$$
$$\left[ \begin{array}{ccc|c} 2 & -1 & 6 & 0 \\ 2 & -1 & 6 & 0 \\ 2 & -1 & 6 & 0 \end{array} \right] \sim \left[ \begin{array}{ccc|c} 2 & -1 & 6 & 0 \\ 0 & 0 & 0 & 0 \\ 0 & 0 & 0 & 0 \end{array} \right]$$
$$2x_1 - x_2 + 6x_3 = 0 \implies x_1 = \frac{1}{2}x_2 - 3x_3$$
Hence the general solution for the eigenvector $\mathbf{x}$ is:
$$\mathbf{x} = \begin{bmatrix} x_1 \ x_2 \\ x_3 \end{bmatrix} = \begin{bmatrix} \frac{1}{2}x_2 - 3x_3 \\ x_2 \\ x_3 \end{bmatrix} = x_2 \begin{bmatrix} 1/2 \\ 1 \\ 0 \end{bmatrix} + x_3 \begin{bmatrix} -3 \\ 0 \\ 1 \end{bmatrix}$$
The eigenspace is a two-dimensional subspace (a plane) of $\mathbb{R}^3$. A basis is
$$\left\{\begin{bmatrix} 1/2 \\ 1 \\ 0 \end{bmatrix}, \quad \begin{bmatrix} -3 \\ 0 \\ 1 \end{bmatrix}\right\}$$

\subsection{Finding eigenvalues}
We have the foundational equation
$$A\mathbf{x} = \lambda\mathbf{x}$$
And we want a nontrivial solution $\mathbf{x}$ to
$$(A - \lambda I)\mathbf{x} = \mathbf{0}$$
By Invertible Matrix Theorem, the matrix equation has nontrivial solutions only if
$$\det(A - \lambda I) = 0$$
We can then derive the characteristic polynomial where the only unknown is $\lambda$. The roots of this polynomial are the eigenvalues.
\paragraph{Example} Find eigenvalues for matrix $A = \begin{bmatrix} 2 & 3 \\ 3 & -6 \end{bmatrix}$.\\
$$A - \lambda I = \begin{bmatrix} 2-\lambda & 3 \\ 3 & -6-\lambda \end{bmatrix}$$
For $2\times 2$ matrix, determinant is $ad-bc$.
\begin{align*}
    \det(A - \lambda I) &= (2-\lambda)(-6-\lambda) - (3)(3)\\
    &= \lambda^2 + 4\lambda - 21\\
    &= (\lambda - 3)(\lambda + 7)
\end{align*}
Setting the characteristic polynomial to zero, we get the eigenvalues $\lambda = 3$ and $\lambda = -7$.

% \subsubsection{Algebraic vs geometric multiplicity}
% Algebraic multiplicity is the number of times a specific eigenvalue appears as a root of the characteristic polynomial.\\
% For example, if a characteristic equation factors to $(\lambda - 5)^2(\lambda - 3)(\lambda - 1) = 0$, the root $\lambda = 5$ occurs twice.
% Therefore, the eigenvalue 5 has an algebraic multiplicity of 2.\\
% \\
% Just because an eigenvalue has an algebraic multiplicity of 2 (it shows up twice in the polynomial), it does not guarantee you will get 2 basis vectors (a geometric multiplicity of 2) when you go to find its eigenspace.
% Always perform row reduction on $(A - \lambda I)\mathbf{x} = \mathbf{0}$ to count your free variables and find the true geometric dimension!

\subsubsection{Special case - triangular matrices}
The eigenvalues of a triangular matrix (upper, lower or pure diagonal) are simply the entries on its main diagonal.\\
\textbf{Example} For lower triangular matrix,
$$B = \begin{bmatrix} \mathbf{2} & 0 & 0 \\ 4 & \mathbf{2} & 0 \\ -1 & 6 & \mathbf{9} \end{bmatrix}$$
The eigenvalues are $\lambda_1 = 2$ and $\lambda_2 = 9$.\\
\textit{Proof} $$ \det \begin{bmatrix} 2-\lambda & 0 & 0 \\ 4 & 2-\lambda & 0 \\ -1 & 6 & 9-\lambda \end{bmatrix} = (2-\lambda)^2\:(9-\lambda) = 0$$
Note that $\lambda = 2$ has an algebraic multiplicity of 2. To find the actual number of independent eigenvectors (geometric multiplicity), you would still need to row reduce $(B - 2I)\mathbf{x} = \mathbf{0}$.\\
\\
\textbf{Example} For pure diagonal matrix,
$$C = \begin{bmatrix} \mathbf{4} & 0 \\ 0 & \mathbf{-1} \end{bmatrix}$$
The eigenvalues are $\lambda_1 = 4$ and $\lambda_2 = -1$.\\
In addition for pure diagonal matrices, the standard basis vectors are automatically their eigenvectors.\\
For $\lambda = 4$, the eigenvector is $\mathbf{v_1} = \begin{bmatrix} 1 \\ 0 \end{bmatrix}$ and For $\lambda = -1$, the eigenvector is $\mathbf{v_2} = \begin{bmatrix} 0 \\ 1 \end{bmatrix}$.

\subsection{Diagonalisation}
Two $n \times n$ matrices $A$ and $B$ are similar if there exists some invertible matrix $P$ such that $B = P^{-1}AP$.\\
Similar matrices represent the same linear transformation, just viewed from different coordinate systems (bases).\\
Importantly, similar matrices share the same
\begin{enumerate}
    \item Determinant
    \item Invertibility
    \item Rank (dimension of column space)
    \item Nullity (dimension of null space)
    \item Trace
    \item Characteristic polynomial
    \item Eigenvalues
    \item Eigenspace dimension. The eigenspace of $A$ corresponding to $\lambda$ and the eigenspace of $P^{-1}AP$ corresponding to $\lambda$ have the same dimension.
\end{enumerate}
In the context of eigenvectors, we want to find a $P$ that transforms $A$ into a diagonal matrix $D$ (all non-zero elements are on the main diagonal).\\
Why diagonalise?
\begin{enumerate}
    \item Eigenvalues of $D$ = the diagonal elements = Eigenvalues of $A$ (by similarity).
    \item $\det(D)$ = product of diagonal entries.
    \item Rank of $D$ = number of non-zero entries on the main diagonal.
    \item Easy multiplication:
    \begin{enumerate}
        \item $DA$: rows of $A$ are multiplied by diagonal elements of $D$.
        \item $AD$: columns of $A$ are multiplied by diagonal elements of $D$.
    \end{enumerate}
    \item $D^k$ just raises each diagonal entry to $k$-th power.
    \item Inverse of $D$ is reciprocal of diagonal elements, i.e. for $A = \begin{bmatrix}a & 0 \\ 0 & b\end{bmatrix}$, $A^{-1} = \begin{bmatrix}\frac{1}{a} & 0 \\ 0 & \frac{1}{b}\end{bmatrix}$.
\end{enumerate}

\subsubsection{Diagonalisation Theorem}
\paragraph{Theorem} An $n \times n$ matrix $A$ is diagonalizable if and only if it has $n$ linearly independent eigenvectors $\mathbf{v}$. If it does, you have enough eigenvectors to form a complete basis for $\mathbb{R}^n$.
$A$ can then be rewriten as $A = PDP^{-1}$.
\begin{enumerate}
    \item $P$ is constructed using the $n$ linearly independent eigenvectors $\mathbf{v}$ as the columns of the matrix.
    \item $D$ is the diagonal matrix built by placing the corresponding eigenvalues on the main diagonal.
    \item The order matters. If $\lambda_1$ is in the first column of $D$, its corresponding eigenvector $\mathbf{v}_1$ must be the first column of $P$.
\end{enumerate}

\subsubsection{$n$ Distinct Eigenvalues - sufficient condition for diagonalisablility}
\paragraph{Theorem} An $n \times n$ matrix $A$ with $n$ distinct eigenvalues is guaranteed to be diagonalizable.\\
This is because:
\\
\textbf{Theorem} If $\mathbf{v}_1,\ldots,\mathbf{v}_r$ are eigenvectors that correspond to distinct eigenvalues $\lambda_1,\ldots,\lambda_r$ of an $n\times n$ matrix $A$,
then the set $\{\mathbf{v}_1,\ldots,\mathbf{v}_r\}$ is linearly independent.\\
Therefore, $n$ unique eigenvalues gives $n$ linearly independent eigenvectors $\mathbf{v}$, fulfilling the condition for the diagonalisation theorem.\\

\subsubsection{Repeated eigenvalues - algebraic vs geometric multiplicity}
A matrix with repeated eigenvalues can still be diagonalisable.\\
First define terms:
\begin{itemize}
    \item Algebraic multiplicity: The number of times an eigenvalue $\lambda_k$ appears as a root in the characteristic equation.
    \item Geometric multiplicity: The actual dimension of the eigenspace (i.e., the number of linearly independent eigenvectors $\mathbf{v}$) corresponding to some eigenvalue $\lambda_k$.\\
    It is calculated as the nullity of ($A - \lambda_k I$), or the number of free variables (columns with no pivots) after row reducing the augmented matrix of $(A - \lambda_k I) \mathbf{v} = \mathbf{0}$
\end{itemize}
For any eigenvalue, 
$$1 \leq \text{Geometric multiplicity} \leq \text{Algebraic multiplicity}$$
For a matrix to be diagonalisable, its geometric multiplicity must exactly equal its algebraic multiplicity for every single eigenvalue.

\paragraph{Example - diagonalisable w/ repeated roots} Given the matrix
$$A = \begin{bmatrix} 1 & 3 & 3 \\ -3 & -5 & -3 \\ 3 & 3 & 1 \end{bmatrix}$$
Skipping the mechanics of finding the determinant,
$$0 = \det(A - \lambda I) = -\lambda^3 - 3\lambda^2 + 4 = -(\lambda - 1)(\lambda + 2)^2$$
The eigenvalue -2 has multiplicity of 2.\\
For $\lambda = 1$: Solving $(A - 1I)\mathbf{x} = \mathbf{0}$ yields one basis vector: $\mathbf{v_1} = \begin{bmatrix} 1 \\ -1 \\ 1 \end{bmatrix}$.\\
For $\lambda = -2$: Solving $(A - (-2)I)\mathbf{x} = \mathbf{0}$ yields two free variables, which allows us to extract two linearly independent basis vectors: $\mathbf{v_2} = \begin{bmatrix} -1 \\ 1 \\ 0 \end{bmatrix} \quad \text{and} \quad \mathbf{v_3} = \begin{bmatrix} -1 \\ 0 \\ 1 \end{bmatrix}$.\\
So we get 3 indepedent eigenvectors, forming a basis for $\mathbf{R}^3$. $A$ is diagonalisable.\\
To diagonalise $A$,
$$P = \begin{bmatrix} \mathbf{v_1} & \mathbf{v_2} & \mathbf{v_3} \end{bmatrix} = \begin{bmatrix} 1 & -1 & -1 \\ -1 & 1 & 0 \\ 1 & 0 & 1 \end{bmatrix}$$
$$D = \begin{bmatrix} 1 & 0 & 0 \\ 0 & -2 & 0 \\ 0 & 0 & -2 \end{bmatrix}$$

\paragraph{Example - Defective matrix} Given the matrix
$$A = \begin{bmatrix} 2 & 4 & 3 \\ -4 & -6 & -3 \\ 3 & 3 & 1 \end{bmatrix}$$
Again skipping the mechanics, the characteristic equation is:
$$0 = \det(A - \lambda I) = -(\lambda - 1)(\lambda + 2)^2$$
For $\lambda = 1$: The basis yields one vector: $\mathbf{v_1} = \begin{bmatrix} 1 \\ -1 \\ 1 \end{bmatrix}$.\\
For $\lambda = -2$: When we solve $(A - (-2)I)\mathbf{x} = \mathbf{0}$ for this new matrix, row reduction only yields one free variable.
This means the eigenspace is only one-dimensional, giving us a single basis vector: $\mathbf{v_2} = \begin{bmatrix} -1 \\ 1 \\ 0 \end{bmatrix}$.\\
Geometric multiplicity of $\lambda = -2$ is 1 $\leq$ its the algebraic multiplicity of 2.\\
Hence $A$ is not diagonalisable (also called a defective matrix).

\subsection{Powers of $A$ - advantage of diagonal matrices}
For a diagonal matrix $D$ and any power $k$,
$$D^k = \begin{bmatrix} 5^k & 0 \\ 0 & 3^k \end{bmatrix}$$
We can derive that for a diagonalisable matrix $A$ and any power $k \geq 1$,
$$A^k = PD^kP^{-1}$$
\paragraph{Example} Compute $A^8$ for the matrix $A = \begin{bmatrix} 4 & -3 \\ 2 & -1 \end{bmatrix}$.\\
Find eigenvalues:
$$\det(A - \lambda I) = \lambda^2 - 3\lambda + 2 = (\lambda - 2)(\lambda - 1) = 0$$
The eigenvalues are $\lambda = 2$ and $\lambda = 1$.\\
Then finding the null space for each $\lambda$, we get the basis vectors v$_1 = \begin{bmatrix} 3 \\ 2 \end{bmatrix}$ and v$_2 = \begin{bmatrix} 1 \\ 1 \end{bmatrix}$.\\
Constructing $P$ and $D$:
$$P = \begin{bmatrix} 3 & 1 \\ 2 & 1 \end{bmatrix}, \quad D = \begin{bmatrix} 2 & 0 \\ 0 & 1 \end{bmatrix}$$
Also,
$$P^{-1} = \begin{bmatrix} 1 & -1 \\ -2 & 3 \end{bmatrix}$$
So $A^8 = PD^8P^{-1}$.
By simple matrix multiplication we get $\begin{bmatrix} 766 & -765 \\ 510 & -509 \end{bmatrix}$.

\subsection{Linear transformation, Eigenbasis}
Applying a linear transformation to a vector $\mathbf{x}$ can be represented as
$$\mathbf{y} = A\mathbf{x}$$
The problem is that in the standard basis, the matrix $A$ is a "coupled" system. When you multiply $A\mathbf{x}$, the rows of $A$ mix the individual $x_1$, $x_2$, and $x_n$ components together.
If you need to apply this transformation multiple times (like $A^{100}$), this coupled mixing requires an enormous amount of computational operations.\\
Instead, we can change the basis to the Eigenbasis, i.e. a basis for a vector space consisting entirely of eigenvectors of a linear transformation or matrix.\\
We use the fact that $A = PDP^{-1}$ to rewrite the linear transformation:
$$A\mathbf{x} = PDP^{-1}\mathbf{x}$$
\begin{enumerate}
    \item $P^{-1}\mathbf{x}$ (Change of basis). $P^{-1}$ is the inverse of the matrix containing the eigenvectors.\\
    This calculates $[\mathbf{x}]_\mathcal{B}$, which are the coordinates of the vector relative to the Eigenbasis. (Recall the change of basis formula $\mathbf{x} = P_\mathcal{B}[\mathbf{x}]_\mathcal{B}$).
    \item $D(P^{-1}\mathbf{x})$ (Transformation). Mulitplying the new coordinates by $D$, the diagonal matrix of eigenvalues.\\
    Because the vector is now defined in terms of eigenvectors, this transformation is strictly a scaling operation. 
    By definition, a transformation only scales eigenvectors by their eigenvalues ($A\mathbf{v} = \lambda\mathbf{v}$).
    The matrix $D$ executes this decoupled scaling by multipliying the first coordinate by $\lambda_1$, the second by $\lambda_2$ etc.
    \item $P(DP^{-1}\mathbf{x})$ (Change back the basis). After the transformation, convert back to standard basis using $P$.
\end{enumerate}
Overall, performing linear algebra operations on a diagonal matrix is supposedly more efficient than operating on a dense, coupled matrix.

\end{document}